% ============================================================================
\clearpage
\section*{v4.10.0-r37 Source-Only: A Stretched String as the Directed OP5 Mediator}
\addcontentsline{toc}{section}{v4.10.0-r37 Source-Only: A Stretched String as OP5 Mediator}
\label{sec:v4100-r37-string-mediator}
% ============================================================================

\paragraph{Purpose.}
After v4.10.0-r36 a vacuum background reproduces the numerator \(N_m=\tfrac23(x^2-a_\perp x+a_\perp^2)\) of the OP5 comparison model, but frame, axis and
polarization structure are chosen.  This layer follows the author's choice of the route through the deferred idea I1: a string stretched
between the matter and the tachyon system, whose tension depends on the Higgs radius.  It is an exploration with a method fixed before the
computation and is not blind; the tensor structure was sketched analytically when the method was written.  Earlier layers are not rewritten.

\noindent\textbf{Source-only status.}
This v4.10.0-r37 layer is source-only.  It does not build a PDF, does not assign a new DOI,
and does not authorize a public release.  The v5.3.0 Version DOI
\href{https://doi.org/10.5281/zenodo.24434403}{10.5281/zenodo.24434403} and the Concept DOI
\href{https://doi.org/10.5281/zenodo.19210728}{10.5281/zenodo.19210728} are unchanged by it.

\subsection*{v4.10.0-r37 assumptions}
\label{subsec:v4100-r37-assumptions}

(S1) A Nambu--Goto string of length \(L\) is stretched along the axis \(n\) between end points at rest in the two systems, with fixed ends;
static gauge, four dimensions, transverse displacements \(X^i\), \(i\in\{x,y\}\).  (S2) Its tension depends on the Higgs radius,
\(T(h)=T_0(1+\tau\eta/v+\dots)\).  (S3) Long wavelength along the axis, \(k_zL\ll1\); the fundamental mode \(\omega_1=\pi/L\) and its two-quantum states are kept.
With \(X^i=q^i s(\sigma)/\sqrt{T_0}\), \(s=\sqrt{2/L}\sin(\pi\sigma/L)\), one has \(\int s^2=1\), \(\int s'^2=\omega_1^2\), \(\int s=2\sqrt{2L}/\pi\), and the second mode has
\(\int s_2=0\).

\subsection*{v4.10.0-r37 results}
\label{subsec:v4100-r37-results}

\paragraph{Theoremlet v4.10.0-r37.1 (the coupling).}
Expanding \(-\int T(h(X))\sqrt{-\gamma}\) to second order in \(X\) gives the terms linear in \(\eta\)
\[
 -\frac{\tau}{v}\sqrt{T_0}\,c_1\,i\,\mathbf k_\perp\!\cdot\mathbf q\,\eta+\frac{\tau}{2v}\,\eta\,\bigl(k_ik_jq^iq^j+\dot q^2-\omega_1^2q^2\bigr),\qquad c_1=\frac{2\sqrt{2L}}{\pi},
\]
and the second mode does not couple at this order.  \emph{Proof.} Insert \(h(X)=\eta\,e^{i\mathbf k\cdot\mathbf X}\) on the string,
\(\sqrt{-\gamma}=1-\tfrac12\dot X^2+\tfrac12X'^2+\dots\), and integrate over \(\sigma\) with the stated mode integrals.  \(\square\)

\paragraph{Theoremlet v4.10.0-r37.2 (two poles).}
In second order the one-quantum states give a pole at \(\omega_1\) with residue \(a_\perp c_1^2T_0\tau^2/v^2\), and the two-quantum states of the fundamental
mode give a pole at \(2\omega_1\) with amplitudes proportional to \(k_ik_j-2\omega_1^2\delta_{ij}\) and residue
\[
 \frac{\tau^2}{2v^2\omega_1}\bigl(a_\perp^2-4a_\perp\omega_1^2+8\omega_1^4\bigr)\;\overset{x=4\omega_1^2}{\propto}\;2a_\perp^2-2a_\perp x+x^2 .
\]
\emph{Proof.} Matrix elements of \(-L_{\rm int}\) between the ground state and the one- and two-quantum states in the Fock space of two oscillators,
\(q=(a+a^\dagger)/\sqrt{2\omega_1}\), \(\dot q=i\sqrt{\omega_1/2}(a^\dagger-a)\); the residue at energy \(E\) is \(\sum2E|\langle n|H_{\rm int}|0\rangle|^2\).  \(\square\)

\paragraph{Remark v4.10.0-r37.3 (reading).}
Scaled to the \(a_\perp^2\) coefficient of \(N_m\) the residue at \(2\omega_1\) is \(\tfrac23a_\perp^2-\tfrac23a_\perp x+\tfrac13x^2\), so it differs from \(N_m\) by \(-\tfrac13x^2\).
Against the criteria fixed beforehand the string route carries the directed contrast \emph{partially}.  Frame and axis are the rest frame and
the direction of the stretched string, and the polarization structure follows from its transverse oscillations without a choice; both
directed parts of \(N_m\) come out in the exact ratio.  The undirected part is half as large, and an additional pole at \(\omega_1\) with numerator
\(a_\perp\) appears; the mass ratio \(1:2\) of the two poles is a prediction of this picture.  The assumptions S1--S3, in particular the dependence
of the tension on the Higgs radius, are set, not derived; the connector levels of I1 are not tested here.

\subsection*{v4.10.0-r37 frames}
\label{subsec:v4100-r37-frames}

All quantities are given in the rest frame of the end points, stated in the matter frame M.  The tachyon end point enters only through the
fixed boundary condition; no quantity of the tachyon frame is compared.

\subsection*{v4.10.0-r37 register entries}
\label{subsec:v4100-r37-register}

\begin{center}
\small
\begin{tabular}{@{}>{\raggedright\arraybackslash}p{0.34\textwidth}>{\raggedright\arraybackslash}p{0.58\textwidth}@{}}
\toprule
statement & status after v4.10.0-r37\\
\midrule
frame and axis of the directed mediator & from the stretched string, under S1--S3\\
directed parts of \(N_m\) & reproduced in exact ratio at the pole \(2\omega_1\)\\
undirected part, extra pole & \(\tfrac13x^2\) instead of \(\tfrac23x^2\); second pole at \(\omega_1\) with numerator \(a_\perp\)\\
\bottomrule
\end{tabular}
\end{center}

\subsection*{v4.10.0-r37 notes}
\label{subsec:v4100-r37-notes}

\texttt{notes/\allowbreak research/\allowbreak op5\_\allowbreak string\_\allowbreak mediator\_\allowbreak 2026\_\allowbreak 11\_\allowbreak 22/} (method note commit 7ecb9fce,
script \texttt{string\_\allowbreak mediator.wl}, output in \texttt{run/}); I1 in \texttt{notes/\allowbreak research/\allowbreak v4\_\allowbreak 10\_\allowbreak 0\_\allowbreak w1/}.

\subsection*{v4.10.0-r37 ledger and guard update}
\label{subsec:v4100-r37-ledger-guard}

\[
R_{\rm chk}^{v4100\text{-}r37}=(1,1,0)
\]
for (status recorded; computations checked in the cited script; the tension law S2 and the undirected part not derived).

The layer spends no guard:
\[
R_{\rm glob}^{(v4100\text{-}r37,+)}=(0,1,1,1,1),
\qquad
\Gamma_{\rm guard}^{v4100\text{-}r37}=(1,1,1,1,1),
\qquad
\beta_4=1/256\ \text{unspent}.
\]
The layer does not imply exact finite \(C5\) zero closure, exact all-readout terminality, \(C6/T2\), a global photon-clock/RH theorem,
or global all-readout closure.  OP1--OP5, C5 and C6/T2 remain open.
